\documentclass[11pt,a4paper]{article}

% --- Typografie & Sprache ---
\usepackage[utf8]{inputenc}
\usepackage[T1]{fontenc}
\usepackage[ngerman]{babel}
\usepackage{lmodern}
\usepackage{microtype}

% --- Layout & Geometrie ---
\usepackage[a4paper,top=2.5cm,bottom=2.5cm,left=2.3cm,right=2.3cm]{geometry}
\usepackage{titling}
\usepackage{abstract}

% --- Mathematik & Physik ---
\usepackage{amsmath,amssymb,amsfonts,amsthm}
\usepackage{mathtools}
\usepackage{braket}

% --- Tabellen & Grafiken ---
\usepackage{booktabs}
\usepackage{tabularx}
\usepackage{array}
\usepackage{colortbl}
\usepackage{xcolor}

% --- Links & Metadaten ---
\usepackage[colorlinks=true,
            linkcolor=blue!70!black,
            citecolor=blue!70!black,
            urlcolor=blue!70!black,
            pdfauthor={Thomas Hofbauer},
            pdftitle={Mathematischer Freeze und duale spektroskopische Laborpfade des EABC-Modells}]{hyperref}

% Farbcodes für Gate-Status
\definecolor{closedGreen}{RGB}{20, 120, 40}
\definecolor{definedBlue}{RGB}{20, 70, 160}
\definecolor{forecastOrange}{RGB}{180, 80, 10}
\definecolor{openRed}{RGB}{160, 20, 20}

\newcommand{\statusClosed}{\textbf{\textcolor{closedGreen}{GESCHLOSSEN}}}
\newcommand{\statusDefined}{\textbf{\textcolor{definedBlue}{DEFINIERT}}}
\newcommand{\statusForecast}{\textbf{\textcolor{forecastOrange}{FORECAST}}}
\newcommand{\statusOpen}{\textbf{\textcolor{openRed}{OFFENE BRIDGE}}}

% Gestaltung der Abstract-Überschrift
\renewcommand{\abstractname}{\large\bfseries Zusammenfassung}
\setlength{\absleftmargin}{1.0cm}
\setlength{\absrightmargin}{1.0cm}

\begin{document}

% --- Titelkopf ---
\begin{center}
    {\LARGE\bfseries Mathematischer Modell-Freeze und duale spektroskopische Laborpfade der nicht-assoziativen EABC-Phasenraumtopologie\par}
    \vspace{0.8em}
    {\large Technisches Whitepaper \& Formale Spezifikation \texttt{Spec-Bridge-01}\par}
    \vspace{1.2em}
    {\textbf{Thomas Hofbauer}\par}
    {\small Bamberg, Deutschland\par}
    \vspace{0.6em}
    {\small Datum: 27. September 2026\par}
\end{center}

\vspace{0.5em}
\hrule
\vspace{1.0em}

% --- Abstract ---
\begin{abstract}
\noindent
Diese Arbeit konsolidiert die Schnittstelle zwischen der nicht-assoziativen Phasenraum-Topologie des EABC-Modells und der Präzisionsspektroskopie des atomaren Wasserstoffs. Auf der algebraischen Ebene wird der Operator-Block über $\mathbb{Q}$ formal geschlossen: Für ein orthogonales Einheitspaar $e_1, e_2 \in \mathbb{O}$ bildet das normierte Assoziator-Quadrat $P_A = \frac{1}{4}K_{12}$ den exakten orthogonalen Projektor auf das Komplement der Quaternionen-Unteralgebra $\mathbb{H}_{12}$ mit $\operatorname{spec}(K_{12}) = \{0^4, 4^4\}$. Unter der Automorphismengruppe $G_2$ liefert der Gruppen-Twirl erster Ordnung das universelle Moment $\mathcal{M}_1 = 0_{\mathbb{R}1} \oplus \frac{4}{7}I_7$, wodurch der oktonionische Ursprung in linearer Näherung zu einer isotropen Skalarkopplung $\eta_A^{\rm eff} = \frac{4}{7}\eta_A$ entartet. Die genuine Fano-Tensorgeometrie überlebt erst im zweiten Moment $\left.\mathcal{M}_2\right|_{\operatorname{Sym}^2(7)} = \frac{4}{7}\Pi_{\mathbf{1}} \oplus \frac{22}{63}\Pi_{\mathbf{27}}$, während gezeigt wird, dass die diskreten 7 Fano-Projektoren kein sphärisches 2-Design bilden. 

Zur empirischen Überprüfung formulieren wir zwei komplementäre Laborpfade: \textbf{Pfad I} fungiert als quantitativer Präzisionsfilter. Durch Einbindung des Mikrowellen-Lamb-Shifts ($2S_{1/2} - 2P_{1/2}$) wird die Kollinearität im optischen $2S\text{--}1S$-Kanal aufgebrochen ($\kappa_2(F) \sim 10^7$), was einen Modell-Forecast von $\sigma_{\rm forecast}(\eta_A^{\rm eff}) \lesssim 1.2 \times 10^{-13}$ ermöglicht. \textbf{Pfad II} adressiert als geometrischer Falsifikator die interne $\mathbf{27}_{G_2}$-Tensorstruktur im Doppler-freien $1S \to 2S$-Zwei-Photonen-Nulltest bei orthogonaler Polarisation ($\boldsymbol{\epsilon}_1 \perp \boldsymbol{\epsilon}_2$). Die Kopplung wird über einen Clebsch-Gordan-Intertwiner $\Xi : \mathbf{27}_{G_2} \to \operatorname{Sym}_0^2(\mathbb{R}^3)$ mit Hebelarm $\operatorname{Tr}(\Xi \mathcal{M}_2) = \frac{110}{63}$ parameterfrei fixiert. Bei Berücksichtigung des relativistischen $E1\text{--}M1$-Standardmodell-Untergrunds ergibt sich eine Sensitivität von $\eta_A \lesssim 10^{-2}$ (quartische Rate) bzw. $\sim 3 \times 10^{-3}$ (Interferenz bei $45^\circ$). Damit trennt das Dossier strikt zwischen mathematischem Beweis, Parameterschärfe und struktureller Falsifikation.
\end{abstract}

\vspace{0.8em}
\noindent\textbf{Keywords:} Oktonionenalgebra, $G_2$-Automorphismen, Fano-Projektor, $1S\text{--}2S$ Spektroskopie, Lamb-Shift, Fisher-Matrix, Clebsch-Gordan-Intertwiner, 2-Design.

\vspace{1.2em}
\hrule
\vspace{1.5em}

% --- Leitmotiv-Box ---
\begin{center}
\colorbox{gray!10}{
    \parbox{0.95\textwidth}{
        \centering\vspace{0.4em}
        \textbf{Methodisches Leitmotiv der Modellvalidierung:}\\[0.2em]
        \emph{Pfad I misst skalare Stärke (Quantitativer Filter via $\mathcal{M}_1 = \frac{4}{7}I_7$).}\\
        \emph{Pfad II misst interne Struktur (Geometrischer Falsifikator via $\mathcal{M}_2 \supset \mathbf{27}_{G_2}$).}
        \vspace{0.4em}
    }
}
\end{center}

\vspace{1.5em}

% --- Audit-Tabelle ---
\subsection*{Formale Gate-Audit-Checkliste (\texttt{Spec-Bridge-01})}
\vspace{-0.5em}

\begin{table}[h!]
\centering
\small
\begin{tabularx}{\textwidth}{@{} c l l X c @{}}
\toprule
\textbf{Gate} & \textbf{Sektor} & \textbf{Objekt / Theorem} & \textbf{Mathematischer / Physikalischer Befund} & \textbf{Status} \\
\midrule
\texttt{G-01} & Algebra & $P_A = \frac{1}{4}K_{12}$ & Exakter Projektor über $\mathbb{Q}$, $\operatorname{spec}(K_{12}) = \{0^4, 4^4\}$ & \statusClosed \\
\texttt{G-02} & Algebra & $\ker(K_{12}) = \mathbb{H}_{12}$ & Kern ist Quaternionen-Unteralgebra $\operatorname{span}_{\mathbb{Q}}\{1, e_1, e_2, e_1e_2\}$ & \statusClosed \\
\texttt{G-03} & Geometrie & $\operatorname{Stab}_{G_2}(P_A) \cong SO(4)$ & Fixierung der Unteralgebra $\mathbb{H}$ bestimmt $SO(4)$ & \statusClosed \\
\texttt{G-04} & Vakuum & $\mathcal{M}_{\rm vac} \cong G_2/SO(4)$ & 8-dimensionale Mannigfaltigkeit der internen Projektoren & \statusClosed \\
\texttt{G-05} & 1. Moment & $\mathcal{M}_1^{\rm Haar} = 0 \oplus \frac{4}{7}I_7$ & Schursches Lemma auf $\mathbb{R}1 \oplus \operatorname{Im}(\mathbb{O})$; kein $\mathbb{O}$-Signal in $\mathcal{O}(1)$ & \statusClosed \\
\texttt{G-06} & 2. Moment & $\mathcal{M}_2^{\rm Haar}\big|_{\operatorname{Sym}^2} = \frac{4}{7}\Pi_{\mathbf{1}} \oplus \frac{22}{63}\Pi_{\mathbf{27}}$ & Normenschranke $M_2 \le I$ erfüllt; $c_{\mathbf{1}}=\frac{4}{7}, c_{\mathbf{27}}=\frac{22}{63}$ & \statusClosed \\
\texttt{G-07} & Sphärik & Fano-2-Design & 7 Fano-Projektoren spalten auf $\mathbf{27}$ in $\frac{4}{7}$ und $\frac{2}{7}$ auf ($1$-Design, kein $2$-Design) & \statusClosed \\
\texttt{G-08} & Intertwiner & $\Xi : \mathbf{27}_{G_2} \to \operatorname{Sym}_0^2(\mathbb{R}^3)$ & Kanonischer Clebsch-Gordan-Projektor via $\mathfrak{so}(3) \hookrightarrow \mathfrak{g}_2$; $\operatorname{Tr}(\Xi \mathcal{M}_2) = \frac{110}{63}$ & \statusDefined \\
\texttt{G-09} & Design & Kollinearitätsbrechung & $2S_{1/2}\text{--}2P_{1/2}$ Lamb-Shift bricht Rydberg-Entartung ($\partial_{\epsilon_h}\Delta E^{(0)}=0$) & \statusClosed \\
\texttt{G-10} & Pfad I & Fisher-Matrix $F$ & Forecast-Sensitivität $\sigma_{\rm forecast}(\eta_A^{\rm eff}) \lesssim 1.2 \times 10^{-13}$ & \statusForecast \\
\texttt{G-11} & Pfad II & Nullkanal $\Gamma_\perp$ & Nichtrelativistischer E1E1-Limes $\mathcal{A}_\perp=0$; Rate $\Gamma_\perp \propto \eta_A^4$ & \statusDefined \\
\texttt{G-12} & Kopplung & $V_{\rm coupling}(\mathbf{r}, \mathbf{p})$ & Dynamischer Ursprung des mikroskopischen Nahfeld-Operators & \statusOpen \\
\bottomrule
\end{tabularx}
\end{table}

\vfill
\noindent\footnotesize\textit{Hinweis zur Verifikation:} Die Sektoren \texttt{G-01} bis \texttt{G-07} sind über $\mathbb{Q}$ bzw. $\mathbb{Z}/2\mathbb{Z}$ in SageMath 10.5 und Lean 4 formal auditiert und als unveränderlicher algebraischer Kern eingefroren.

\end{document}